\SourceLeaf{340}
also:
\[
\begin{array}{rl}
(4,1) : (4,\text{III}) & = (3,1) : (\text{I}, \text{III}) \\
(4,1) : (4,\text{II}) & = (2,1) : (\text{I}, \text{II}) \\
(4,3) : (4,\text{II}) & = (2,3) : (\text{III}, \text{II}).
\end{array}
\]
Die beiden letzten Gleichungen geben durch Division:
\begin{equation}
(4,1) : (4,3) = \frac{(2,1)}{(2,3)} : \frac{(\text{I}, \text{II})}{(\text{III}, \text{II})}
= \frac{(2,1)}{(3,2)} : \frac{(\text{I}, \text{II})}{(\text{II}, \text{III})}.
\tag{1}
\end{equation}
Fügt man zum System der drei Gleichungen noch hinzu:
\[
(4,2) : (4,\text{III}) = (3,2) : (\text{II}, \text{III}),
\]
so hat man
\begin{equation}
(4,1) : (4,2) = \frac{(3,1)}{(3,2)} : \frac{(\text{I}, \text{III})}{(\text{II}, \text{III})}
\tag{2}
\end{equation}
und ebenso:
\begin{equation}
(4,2) : (4,3) = \frac{(2,1)}{(3,1)} : \frac{(\text{I}, \text{II})}{(\text{I}, \text{III})}.
\tag{3}
\end{equation}
Aus (1) folgt:
\[
\frac{(\text{I}, \text{II})}{(\text{II}, \text{III})} = \frac{(2,1) . (4,3)}{(4,1) . (3,2)},
\]
aus (2) folgt:
\[
\frac{(\text{I}, \text{III})}{(\text{II}, \text{III})} = \frac{(3,1) . (4,2)}{(4,1) . (3,2)}.
\]
Demnach ist:
\[
(\text{I}, \text{II}) : (\text{II}, \text{III}) : (\text{III}, \text{I}) = (2,1) . (4,3) :
(4,1) . (3,2) : (4,2) . (3,1).
\]
Bezeichnen wir die Winkel wie vorher, und setzen:
\[
\sphericalangle\ \text{I}, \text{II}, \text{III} = [\text{I}, \text{III}],
\]
so ist:
\[
(\text{I}, \text{II}) : (\text{II}, \text{III}) : (\text{III}, \text{I}) = \sin[\text{I}, \text{II}] :
\sin[\text{II}, \text{III}] : \sin[\text{III}, \text{I}].
\]
Nun ist:
\[
\begin{array}{rl}
{[\text{I}, \text{II}]} & = [\text{I}, 4] + [4, \text{II}] = \alpha + \delta \\
{[\text{II}, \text{III}]} & = [\text{II}, 4] + [4, \text{III}] = \gamma + \zeta \\
{[\text{III}, \text{I}]} & = [\text{III}, 4] + [4, \text{I}] = \beta + \varepsilon,
\end{array}
\]
also:
\[
\frac{\sin(\alpha+\delta)}{(2,1) . (4,3)} = \frac{\sin(\gamma+\zeta)}{(4,1) . (3,2)} =
\frac{\sin(\beta+\varepsilon)}{(4,2) . (3,1)},
\]
welches mit der frühern Gleichung übereinstimmt.

Ähnlich stellt sich die Aufgabe, liegt 4 ausserhalb des Dreiecks 1, 2, 3. Man nimmt dann beliebig
Eine der vier Ecken, z.\,B. wieder 4, zieht (4,1), (4,2), (4,3) und bestimmt ähnlich die Punkte
I, II, III und hierauf das Dreieck (I), (II), (III), indem man wieder die Relationen zwischen
Seiten und Winkeln aufsucht.
